\documentclass{article}

\usepackage{arxiv}

\usepackage[utf8]{inputenc} % allow utf-8 input
\usepackage[T1]{fontenc}    % use 8-bit T1 fonts
\usepackage[hidelinks]{hyperref}       % hyperlinks
\usepackage{url}            % simple URL typesetting
\usepackage{booktabs}       % professional-quality tables
\usepackage{amsmath,amssymb,amsthm}
\usepackage{amsfonts}       % blackboard math symbols
\usepackage{nicefrac}       % compact symbols for 1/2, etc.
\usepackage{makecell}
\usepackage{microtype}      % microtypography
\usepackage{mathrsfs}
\usepackage{float}
\usepackage{graphicx}
\usepackage{doi}
\usepackage{acronym}
\usepackage{listings}
\usepackage{tabularx}
\usepackage{tikz}
\usepackage[dvipsnames]{xcolor}

\usetikzlibrary{trees}

\def\checkmark{\tikz\fill[scale=0.4](0,.35) -- (.25,0) -- (1,.7) -- (.25,.15) -- cycle;}
\def\cross{\tikz\draw[scale=0.3, black, line width=0.3mm](0,0) -- (1,1) -- (0.5,0.5) -- (0,1) -- (1,0) -- (0.5,0.5);}

\newacro{abm}[ABM]{Agent-Based Model}
\newacro{cabm}[CABM]{Cellular Agent-Based Model}
\newacro{ca}[CA]{Cellular Automaton}
\newacro{ib}[IB]{Individual-Based}
\newacroplural{ca}[CA]{Cellular Automata}
\newacro{ecoli}[\textit{E.coli}]{\textit{Escherichia coli}}
\newacro{bsubtilis}[\textit{B.subtilis}]{\textit{Bacillus subtilis}}
\newacro{spombe}[\textit{S.pombe}]{\textit{Schizosaccharomyces pombe}}
\newacro{paeruginosa}[\textit{P.aeruginosa}]{\textit{Pseudomonas aeruginosa}}

\newcommand{\todo}[1]{\colorbox{WildStrawberry}{\textcolor{white}{#1}}}

\newcommand{\R}{\mathbb{R}}

\title{Mechanical Models of Rod-Shaped Bacteria - From Single-Cell Essays to Collective Phenomena}

%\date{September 9, 1985}	% Here you can change the date presented in the paper title
%\date{} 					% Or removing it

\author{
    \href{https://orcid.org/0009-0001-0613-7978}{
        \includegraphics[scale=0.06]{orcid.pdf}
        \hspace{1mm}Jonas Pleyer
    }
    \thanks{
        \href{https://jonas.pleyer.org}{jonas.pleyer.org},
        \href{https://cellular-raza.com}{cellular-raza.com}
    }\\
	Freiburg Center for Data-Analysis and Modeling\\
	University of Freiburg\\
	\texttt{jonas.pleyer@fdm.uni-freiburg.de} \\
	%% examples of more authors
	\And
	\href{https://orcid.org/0000-0002-6371-4495}{
        \includegraphics[scale=0.06]{orcid.pdf}
        \hspace{1mm}Christian Fleck
    }\\
	Freiburg Center for Data-Analysis and Modeling\\
	University of Freiburg
}

% Uncomment to remove the date
%\date{}

% Uncomment to override  the `A preprint' in the header
\renewcommand{\headeright}{Preprint}
%\renewcommand{\undertitle}{Technical Report}
\renewcommand{\shorttitle}{Mathematics of rod-shaped bacteria}

\usepackage{enumitem}
\setlist{nolistsep}

%%% Add PDF metadata to help others organize their library
%%% Once the PDF is generated, you can check the metadata with
%%% $ pdfinfo template.pdf
\hypersetup{
pdftitle={Mathematics of rod-shaped bacteria},
pdfsubject={q-bio.NC, q-bio.QM},
pdfauthor={Jonas Pleyer, Christian Fleck},
pdfkeywords={},
}

% Define definition, example, lemma, proof and theorem.
\newtheorem{definition}{Definition}[section]
\newtheorem{example}[definition]{Example}
\newtheorem{lemma}[definition]{Lemma}
\newtheorem{corollary}[definition]{Corollary}
\newtheorem{theorem}[definition]{Theorem}

% Change numbering of equations
% \numberwithin{equation}{section}

% MAKE TITLES IN THEOREMS BOLD
\makeatletter
\def\th@plain{%
  \thm@notefont{}% same as heading font
  \itshape % body font
}
\def\th@definition{%
  \thm@notefont{}% same as heading font
  \normalfont % body font
}
\makeatother

\begin{document}
\maketitle

% TABLE OF CONTENTS
% Remove this before submission

%###################################################################################################
\begin{abstract}
\end{abstract}

% keywords can be removed
\keywords{Single-Cell \and Mechanics \and Rod-Shaped \and Bacteria \and Mathematics}

\textbf{POSSIBLE TITLES}
\begin{itemize}
    \item Mechanical Models of Rod-Shaped Bacteria - From Single-Cell Modeling to Collective Phenomena
    \item Mathematical Modeling of Rod-Shaped Bacteria - From Single-Cell Mechanics to Collective Phenomena
    \item Mathematical Models of Rod-Shaped Bacteria - From Single-Cell Analyses to Collective Phenomena
\end{itemize}

\tableofcontents

%###################################################################################################
\section{Introduction}

\begin{itemize}
    \item Interested in all variations of rod-shaped bacteria; stiff, flexible, spiral, etc.
    \item Einordnung in andere Bakterienarten
    \item Rods are "more fundamtenal" than spheres (cocci) (which citation?)
\end{itemize}

\section{Biological Principles}
\subsection{How Bacteria grow and maintain their Shape}

The last decades have fundamentally challenged and changed our understanding of how bacteria retain
their shape.
Rod-shaped bacteria can grow by extending their circular part, or by inserting new material
at the tip of the rod.
Species such as \ac{ecoli} and ac{bsubtilis}~\cite{Errington2020} fall under the first category
while \ac{spombe} represents the latter.
The material is inserted in small bursts on the nanoscale in forms of patches, bands or
hoops~\cite{DePedro2003}.
Despite their differences in wall-size, \ac{bsubtilis} and \ac{ecoli} follow growth rules which are
comparable with respect to the extension of the rod \cite{Chang2014}.

It was long believed that bacteria lack cytoskeletal filaments but with the end of the last
century, it was discovered that the MreB protein fills this role since it "[..] forms an actin-like
cytoskeleton in bacteria [..]" \cite{Erickson2001} (supplement: [p. 1]).
It can polymerize and form filaments which behave similar to actin microfilaments \cite{Dersch2020}.
Studies of its crystal structure further showed the similarity to actin~\cite{vandenEnt2001} and the
MreB, MreC and MreD proteins were identified as homologues of actin \cite{Lowe2017_lj}.\
\cite{Jones2001} studied the species of different bacterial sub-kingdoms which have the MreB gene.
They concluded: "Many of the organisms are rod shaped, similar to B. subtilis and E. coli.
Organisms with more complex shapes, including curved, filamentous, and helical bacteria, were all
abundantly represented." \cite{Jones2001} (supplement: [p. 7]).
Furthermore \cite{Wachi1987} found that mutants of \ac{ecoli} which create defective MreB proteins will be
spherical instead of their natural rod-shaped form.
Our colleagues from Freiburg have validated many of these observations and further showed that MreB
travels along the axis of elongation with a velocity that depends on the size of the filament
\cite{Olshausen2013}.

\begin{itemize}
    \item z-Ring formation (citations and text)
    \item \cite{Koch1995} TODO "A physical basis for the precise location of the division site of rod-shaped bacteria: the Central Stress Model Free"
    \item \cite{Bramkamp2009} TODO "Division site selection in rod-shaped bacteria"
\end{itemize}

\begin{figure}[H]
    \centering
    \includegraphics[width=0.3\textwidth]{example-image-a}
    \includegraphics[width=0.3\textwidth]{example-image-b}
    \includegraphics[width=0.3\textwidth]{example-image-c}
    \caption{TODO:
        (A) Insert material along Rod
        (B) Insert material at tips
        (C) Z-Ring formation and division process
    }
\end{figure}

This raises the question whether a rod shape serves any particular purpose.
It can be argued that the earliest cells consisted of "[..] rods and filaments with cocci being
derived and degenerate forms." \cite{Young2006}.
The importance of shapes on the function of cells is highlighted by~\cite{Takeuchi2005} who grew
\ac{ecoli} in very narrow microchambers and thus permanently altered their shape.
This resulted in modified motion of the cells, some unchanged (helical and short crescent cells),
some dysfunctional (coiled cells).
It is important to note that this experiment controlled only the shape of bacteria without altering
any additional biochemical or genetic functions.

\paragraph{Bending Stiffness}

As discussed previously, the major contributor which allows rod-shaped bacteria to grow an keep
their shape is MreB.
If left alone, the rod will be straight but external effects can alter the morphology
permanently~\cite{Takeuchi2005}.
Multiple efforts have been made in recent years to measure and analyze the bending properties of
individual bacteria.
\cite{Amir2014_2} designed an experiment similar to \cite{Wang2010} but instead of optically trapped
beads, a flow was used to apply a uniform force on the rod.
Their method is non-invasive and allows for larger force scales.
Cell-division was suppressed with SulA, thus altering the cellular biochemistry.
Since the cells grow continuously, a straight shape was recovered in every case.
With these methods, they demonstrated that the bacteria "[..] exhibit two fundamentally different
modes of deformations and recover from them" \cite{Amir2014_2}.
These modes are elastic in which the cell recovers immediately and plastic which requires growth to
restore its morphology.
Short pulses of forces resulted in elastic deformations from which the cell recovered immediately.
Plastic changes in morphology were only possible under longer periods of persistent force exertion.
Figure X(B) compares the angle of deflection before and after the snap-back
which occurs after turning off the flow.
The linear fit of 66 experimental results under varying conditions shows that both modes contribute
to the behaviour of the cells.

\begin{figure}[H]
    \centering
    \includegraphics[width=0.45\textwidth]{example-image-a}
    \includegraphics[width=0.45\textwidth]{example-image-b}
    \caption{TODO:
        (A) Cell-cell Interactions
        (B) Elastic and Plastic Deformations of the Rod
    }
\end{figure}

\subsection{Physical Interaction between individual Bacteria and Surfaces}

A colony of bacteria is known to form biofilms~\cite{Dunne2002}.
This strategy is is beneficial for the survival of the collective as it allows it to harness growth
factors and migrate when necessary and also plays an important role in bacterial infections in the
biomedical sector \cite{Ong1999}.
The formation of biofilms is only possible due to the attractive nature of bacterial interactions
towards each other and surfaces \cite{Berne2018}.
To analyze bacterial adhesion to surfaces, \cite{vanLoosdrecht1989} investigated the interactions of
various bacteria with negatively charged polystyrene.
They found that the interaction is characterized by a low Gibbs energy of
$2-3kT$ per cell which can be described by a secondary minimum of the DLVO-theory
\cite{Derjaguin1993,Verwey1947} which is relevant for distances $>=1nm$.
At these distances, the DLVO potential acts as a generic adhesive potential.
Particles which stay in the secondary minimum are held together but even weak interactions are
enough to redisperse them which means that this state is reversible.
When approaching closer distances, more intricate mechanisms come into play in which the
heterogeneity of the materials in question needs to be considered.
\cite{Hori2010} described bacterial adhesion as a "two-phase process including an initial, instantaneous
and reversible physical phase (phase one) and a time-dependent and irreversible molecular and
cellular phase (phase two)".

All these considerations become even more intricate when not considering interactions of bacteria
with surfaces but bacteria with each other.
Gram-negative bacteria (eg. \ac{ecoli}) express outer membrane proteins (adhesins) which allow them
to attach to various hosts \cite{Vaca2019,Beachey1981}.
In addition, the flagella can play an important role in the direct adhesion between cells
\cite{Haiko2013}.
Thus it is not correct to blindly apply the results for adhesion to surfaces to the interactions
between bacteria.

The overall attractive nature of the bacterial forces can be seen in the stress response to positive
and negative compression of a colony in a confined space.
\cite{Trejo2013} studied the formation of macroscopic wrinkles when growing bacteria in a confined
space.
They were able to model the strain and stress of the collection of bacteria which requires that
cells are connected by a (to first order) spring-like force.
The developing wrinkles could then be described by a buckling instability.

\cite{Duvernoy2018} used laser techniques and force microscopy to investigate adhesion of individual
\ac{ecoli} and \ac{paeruginosa} cells.
They found that the adhesion forces of the rods are polar, resulting in mechanical tension which in
turn determines daughter cell arrangement.
They measured the size of the grown bacterial colonies, fitted ellipses to their shapes and found
that adhesion and polar adhesion in particular resulted in a greater difference between the long and
short axis of the fitted ellipse and thus a more oval shape rather than circular (WT).
Furthermore, they measured the critical value $N_(2D\/3D)$ where the colony grows into a second
layer on top of the base layer.
It was shown by \cite{Grant2014} Figure X B), that the point of transition
depends on the rigidity of the surrounding gel and \cite{Duvernoy2018} extended this analysis,
showing that the critical value $N_(2D\/3D)$ depends on the adhesive strength and polarity of the
interaction.
\cite{Grant2014} described the growing bacterial colony as a disc with radius $R$ which is pressed
into the agarose slab.
The authors created a purposely-designed simulation in `C++` in order to model the effects of the
gel on the collection of bacteria.
They approximated individual bacteria as multiple overlapping spheres which were linked by nonlinear
springs such that their dynamics correspond to the Euler-Bernoulli dynamic beam theory
\cite{HAN1999}.
The physical interaction potential between multiple bacteria and the surface was of the form

\begin{equation}
    F(\epsilon) = C E_b (d/2)^(2 - \beta) \epsilon^\beta
\end{equation}

where only repulsive effects were considered and no attractive forces which is justified by the
already dense packing of the bacteria in the experiment.
The value $F$ is the magnitude of the force, $E_b$ is the effective Young modulus of the whole
bacterium, $epsilon$ is the distance of overlap between two spheres and $d$ the diameter of the
spheres.
Their model also accounts for friction between cells which is based on Amonton's laws (against the
direction of movement, proportional to the magnitude of the normal force) \cite{Hutchings2021}.
In order to portray growth, new spheres were inserted and cells divided upon reaching a threshold.
In combining all these aspects, the authors have constructed an agent-based model which was able to
produce the stacking behaviour they set out to describe.

\subsection{Emergent Phenomena and Self-Organization}

Although the developing wrinkles mentioned in the preceding section can be characterized as a form
of collective self-organization, they require a confined space and would not be observable
otherwise.
In most cases, rod-shaped bacteria grow in a medium separated enough such that they can freely
extend in space.
\cite{vanGestel2015} explored how \ac{bsubtilis} migrates over a surface by forming multicellular
structures.
The cells organize themselves into "van Gogh bundles" of cells
which are tightly aligned in chains which then form filamentous loops.
This phenomenon occurs at the border of the bacterial colony and the migration is driven by two
phenotypically different cell types \cite{Lpez2010}.
"[The authors] propose that surfactin-producing cells reduce the friction between cells and their
substrate, thereby facilitating matrix-producing cells to form bundles" \cite{vanGestel2015}.
It has been shown that mutants which combine \textit{srfA} and \textit{eps} are able to outperform
the wild-type when it comes to colony expansion \cite{Velicer2009}.
In their final stage, van Gogh bundles consist only of matrix-producing cells although they depend
on surfactin-producing ones during their development.
\cite{vanGestel2015} also constructed an agent-based model to describe the observed effects.
Cells are modeled as straight lines in the shape of a filament by connecting their ends to each
other.
They are initially placed horizontally at the bottom of the simulation domain
In order to update to the next position, three mechanisms are involved: (1) Cell elongation (2) Cell
division (3) Cell turning.
While the first two are rather self-explanatory, in the last step, a new spatial orientation of the
cells is chosen in order to minimize the potential energy

\begin{equation}
    V = (\pi - \alpha_1)^2 + (\pi - \alpha_2)^2
\end{equation}

where $alpha_1$ and $alpha_2$ denote the angles to their respective neighbors.
By choosing the new orientation at random with probability $P=1-e^(k(V_c-V_n))$ (where $V_c,V_n$ are
the current and new potential energies), the parameter $k$ symbolizes the bending rigidity.
Their analysis found that a high bending rigidity can play a more important role than the overall
cell size, i.e. the growth rate.
Concerning the complexity of their model, the authors noted:
"We did not aim to accurately model the biophysical details of the growth of van Gogh bundles
(parameterization of such a model would be impossible), but rather to make a simple phenomenological
model to shape our intuition [..]" \cite{vanGestel2015}.
While the authors correctly assessed that the parametrization of such a model with just one
particular experiment would be very difficult, multiple sources combined which treat individual
aspects of each cellular aspect can yield enough information to meaningfully produce results.
\cite{Dong2022} showed that van Gogh bundles aid in the self-healing of bacterial colonies.
They observed that an artificially introduced cut heals better at lower curvatures, which
corresponds to a higher rigidity of the ensemble being placed more towards the outer region of the
triangular-shaped cut.
This work served as a precursor for a very recent study by \cite{Li2025} in which they picked up on
the Agent-Based model.
Here, they investigated the interplay of nutrient depletion and multiple biofilms on the
capacity to form van Gogh bundles.
Their model used linear growth on the cellular level which is only a first-order approximation to
the experimentally observed exponential growth discussed previously.
Furthermore, no forces were calculated between the bacteria but a "shoving algorithm" (collision
detection) moved the agents around such that collisions were avoided.
Based on the approach of \cite{vanGestel2015} they implemented cell-turning with.
But more significantly, they implemented a reaction-diffusion system

\begin{equation}
    \frac{\partial C_s}{\partial t} = D_n \nabla C_s + r_n
\end{equation}

where $\nabla$ is the gradient, $D_n$ the diffusion coefficient and $r_n$ is the nutrient sink or
source.
More precisely, the nutrient sinks (sources) $r_n$ have to be located at the positions of cells.
If we assume a simplification to a point-like interaction with the external environment, we can
write as

\begin{equation}
    \frac{\partial C_s}{\partial t} = D_n \nabla C_s  + \sum\limits_i u_i r_n (x_i) \delta(x-x_i)
\end{equation}

where the sum ranges over all cells $i=1..$ and each cell might have an individual exchange range
$u_i$ located at position $x_i$ which is denoted by the well-known dirac delta-distribution
$\delta(x-x_i)$.
The authors compared a single-biofilm system with a multi-biofilm system and showed that
motile cells can convert into matrix-producing cells in a nutrient depleted system.
The authors show that spores in the case of a nutrient-depleted system

\section{Mathematical \& Computational Frameworks}

In order to describe the multicellular systems we looked at in the preceding sections, multiple
different aspects of cellular behaviour and their interactions with the surrounding domain need to
be considered and implemented.
Table \ref{table:simulation-aspects-supplement} summarizes these aspects and points to the relevant
experiments.
Aspects (1-5) take place inside the cell, (6,7) correspond to interactions between different
cellular agents and (8,9) to interactions with the domain.
It should be noted explicitly that these simulation aspects can be coupled to each other.
We have already seen such an example in the results of @Takeuchi2005 and @Ursell2014 where the
continued growth modulates the mechanical response of the cell.
Additionally, Aspect (5) is an overarching concept which is relevant for all processes.
Especially for the parameters which facilitate growth, we can assume that there is no direct
inheritance from one generation to the next but only a stochastic distribution of parameters from
which the new value is drawn (see supplement).

\subsection{Single-Cell Modeling Aspects}

\begin{table}[H]
    \centering
    \def\arraystretch{1.3}
    \begin{tabularx}{\textwidth}{c l X}
        &\textbf{Aspect} & \textbf{Description}\\
        \toprule
        &\textbf{(C) Cellular}\\
        \midrule
        (1) & Rod-Shaped Mechanics &
            Rod-shaped bacteria are flexible rods which are able to freely move around (off-lattice
            approach) \cite{Takeuchi2005,Ursell2014,Amir2014_2}.\\
        (2) & Growth &
            Cells grow exponentially by inserting new material either along the circular part of the
            rod or at the tip~\cite{Robert2014,Takeuchi2005}.\\
        (3) & Differentiation &
            \textit{B.subtilis} is known to differentiate into matrix-producing and
            surfactin-producing cells \cite{vanGestel2015,Lpez2010}.\\
        (4) & Division &
            The formation and continuation of van Gogh bundles is driven by cell division
            \cite{vanGestel2015}.
            Bacteria can form multilayers during their growth phase \cite{Duvernoy2018}.\\
        (5) & Variable Parameters &
            Parameters for individual cells are not fixed values but rather taken from a
            distribution \cite{Koutsoumanis2013}.\\
        &\textbf{(CC) Cell-Cell Interactions}\\
        \midrule
        (6) & Adhesion &
            Bacteria adhere to each other at longer distances and attach when in close contact
            \cite{Verwey1947,Trejo2013}.
            The interaction between rods can be polarized \cite{Duvernoy2018}.\\
        (7) & Friction &
            Friction between cells \cite{Grant2014} can be asymmetrical \cite{Doumic2020}.\\
        &\textbf{(DC) Domain-Cell Interactions}\\
        \midrule
        (8) & External Forces &
            Bacteria tend to stick to surfaces \cite{vanLoosdrecht1989}.
            The extracellular gel exerts a force onto the cells \cite{Grant2014}.\\
        (9) & Extracellular Reactions &
            Bacteria can take up nutrients or possible secrete/take up signalling molecules
            \cite{Li2025}.\\
        \bottomrule
    \end{tabularx}
    \label{table:simulation-aspects-supplement}
    \caption{TODO}
\end{table}

\subsection{Computational Modeling Frameworks}

Since the beginning of this decade, multiple tools have emerged which are able to describe living
systems on a cellular individual-based approach in various details.
A good reason for choosing a framework over a purpose-built solution is to follow the Findability,
Accessibility, Interoperability and Reuse) FAIR principles by \cite{Wilkinson2016}.
These criteria have been designed to improve the overall infrastructure surrounding (re)usability of
scholarly data and methods.
In our previous work, we investigated their differences and features and capacity to model
individual behaviour of cells \cite{Pleyer2023}.
Using one of these existing toolkits means that already existing functionality can be used to
develop own models and the produced research results can be fed back to the library for fellow
researchers to reuse.
This begs the question, which of these existing models is able to support the long list of aspects
that we set out to describe with the experimentally gathered evidence.

\paragraph{Many general-purpose Frameworks lack explicit support for Rod-Shaped Bacteria}

Of the most popular frameworks, most do not support rod-shaped bacteria out of the box.
The very popular Cellular Potts Model (CPM) \cite{Graner1992} is mostly applied in 2D and only
represents cells as lattice grid points.
Since CompuCell3D \cite{Swat2012} exclusively builds upon the CPM, it can not model rod-shaped
bacteria.
PhysiCell \cite{Ghaffarizadeh2018} was designed to answer questions surrounding cancer research and
currently only supports spherical agents.
Chaste \cite{Cooper2020} was also designed for cancer research but further targets the heart and
tissues.
Naturally its Agent-Based Model supports cell-centre and vertex-based models but no rod shapes.
Morpheus \cite{Starru2014} also employs the CPM along with other spatial representations such as
vertex-based models or PDEs but has no support for Rod-Shaped bacteria.
Even purpose-built solutions such as BNSim \cite{Wei2013} which specifically targets bacterial
networks, assumes a simplified spherical representation for their bacterial agents.

\paragraph{Varying Support for Rod-Shaped Bacteria}

\begin{figure}[H]
    \centering
    \includegraphics[width=0.3\textwidth]{example-image-a}
    \includegraphics[width=0.3\textwidth]{example-image-b}
    \includegraphics[width=0.3\textwidth]{example-image-c}
    \caption{TODO: Snapshots of Modeling Frameworks}
\end{figure}

As we have seen, many frameworks do not provide existing functionality for rod-shaped bacteria.
However, some others do have various levels of support.
Biocellion \cite{Kang2014} can model agents with cylindrical interaction potentials but does not
model any of the MreB-related bending and rigidity or polar interactions.
They acknowledge this shortcoming: "Mapping a cell to multiple agents is also necessary to
separately model subcellular compartments [..].
However, Biocellion does not yet support this." \cite{Kang2014}
BSim2.0 represents cells as rigid capsular cells made from a cylindrical center part and two
half-spheres, which are placed at the ends of the cylinder to round out the shape.
In order to calculate interactions between cells, possible overlaps are determined and minimized,
thus determining the position values of the next iteration step.
It also accounts for many other phenomena, which are displayed in \textbf{TODO replace}.
BSim does not consider bending forces for individual cells or polar interactions.
The `gro` programming language was designed to simulate the growth of colonies and cell-cell
communication.
Its "physics computation has been optimized for rigid rod-shaped bodies, like E. coli bacteria"
\cite{Gutirrez2017}.
They recognize two types of forces which are acting on the cellular agents:
\textit{Local forces} which are calculated between adjacent bacteria and a \textit{global force}
which pushes bacteria outwards of the colony.
The latter of these is a phenomenological implementation of the observed colony expansion and the
associated central pressure with it.
This assumption may yield incorrect results for sparsely populated cases.
The engine is limited to 2D and does not consider polar interactions or bending of the rods.

\subsection{Purpose-Built Models}

\begin{itemize}
    \item \cite{Abar2017} "Agent Based Modelling and Simulation tools: A review of the state-of-art software"
    \item other generic frameworks exist, which may allow to model the desired functions
    \item \cite{Winkle2017} "Modeling mechanical interactions in growing populations of rod-shaped bacteria"
    \item \cite{Doumic2020} "A purely mechanical model with asymmetric features for early morphogenesis of
  rod-shaped bacteria micro-colony"
    \begin{itemize}
        \item This is a really good paper for referencing
        \item no bending in mathematical model
        \item Compare distributions of "Read-Outs" to data
        \item uses steric force (see also previously \cite{Trejo2013})
    \end{itemize}
    \item \cite{Grant2014} purposely-built model written in `C++`
    \begin{itemize}
        \item describes bacteria as collection of overlapping spheres
        \item spheres are coupled by non-linear springs (Euler-Bernoulli dynamic beam theory);
        This assumption is unfounded; they show however that it does not alter their results in this
        case
        \item only model repulsive forces; no attraction, adhesion
    \end{itemize}
    \item \cite{Cho2007} only 2D, no bending, no parameter estimation; based on work done in 
        \cite{Jnsson2005}
    \item \cite{Storck2014} TODO;
    41 Parameters for various cases;
    only 8 parameters taken from literature values/quantified;
    generated growth rate randomly (normal distribution), but for each growth step; this is
      stochastically equivalent but numerically slightly more intense;
    \item \cite{Volfson2008} TODO; continuum model, equations of nematodynamics \cite{Doi1988-ad}
    \begin{align}
        \partial_t \rho + \partial_z (\rho \nu) &= \alpha \rho\\
        \partial_t q + \nu \partial_z q &= B(1-q^2) \partial_z \nu\\
        \partial_t(\rho \nu) + \nu \partial_z (\rho \nu) &= - \partial_z p - \mu \rho \nu
    \end{align}
    \item \cite{Pleyer2023} TODO "Modeling mechanical interactions in growing populations of rod-shaped bacteria"
    \item \cite{Kong2014} TODO "Swimming motion of rod-shaped magnetotactic bacteria: the effects of shape and growing magnetic moment"
\end{itemize}

\subsection{Parameter Estimation Approaches}

\begin{itemize}
    \item \cite{Storck2014} TODO; 41 Parameters for various cases; only 8 parameters taken from
        literature values/quantified
    \item Highlight lack of estimation of mechanical parameters for Agent-Based Models
\end{itemize}

\section{Discussion}
\section{Conclusion}

\bibliographystyle{IEEEtran}
\bibliography{references}

\renewcommand{\thesection}{}
\renewcommand{\thesubsection}{S\arabic{subsection}}

\section{Supplementary Material}
\subsection{Inheritance of Growth-Relevant Parameters}

We consider a distribution of parameters $rho(x)$ for
which the rate of proliferation of any cell is proportional to the value of the parameter $x$.
The rate of proliferation and thus the rate with which $rho(x)$ changes is given by
$partial_rho(x,t) x = lambda rho(x,t) x$.
This Ordinary Differential Equation (ODE) is trivially solvable with solution

\begin{equation}
    \rho(x,t) = \rho(x,t_0) \exp(\lambda t x)
\end{equation}

For the simple example of a Gaussian distribution
$\rho(x,t_0) = 1/\sqrt{2 pi \sigma^2} \exp(-x^2/(2 \sigma^2))$, we can calculate the expectation
value for the parameter $x$ via

\begin{align}
    E[x(t)]
    &= 1/\sqrt{2 \pi \sigma^2} \int \rho(x,t) x op("dx")\\
    &= 1/\sqrt{2 \pi \sigma^2} \int x \exp{-x^2/(2 \sigma^2} + \lambda t x) op("dx")\\
    % &= 1/sqrt(2 pi sigma^2) integral x exp( -(x - lambda sigma^2 t)^2/(2 sigma^2) - lambda^2 sigma^2 t^2) op("dx")\
    % &= 1/sqrt(2 pi sigma^2) integral x exp( -(x - lambda sigma^2 t)^2/(2 sigma^2) + (lambda^2 sigma^2 t^2) /2) op("dx")\
    % &= 1/sqrt(2 pi sigma^2) integral [- sigma^2 partial/(partial x) +lambda sigma^2 t] exp( -(x - lambda sigma^2 t)^2/(2 sigma^2))exp((lambda^2 sigma^2 t^2) /2) op("dx")\
    % &= 1/sqrt(2 pi sigma^2) integral lambda sigma^2 t exp( -(x - lambda sigma^2 t)^2/(2 sigma^2))exp((lambda^2 sigma^2 t^2) /2) op("dx")\
    &= \lambda \sigma^2 t \exp{(\lambda^2 \sigma^2 t^2)/2}.
\end{align}

We can see that the expected value $E[x(t)]$ keeps increasing faster than regular exponential growth
and indefinitely which is contrary to observation (see Figure \textbf{TODO}) and intuition.
The same qualitative result will emerge for different distributions.
With these considerations, we can assume that the new parameters of freshly divided cells are drawn
from a random distribution which carries no temporal correlation between the mother and daughter
cell. (TODO this should be provable more rigorously)
It has to be stated explicitly that these considerations were done under the assumption that the
proliferation of the cell is affected and proportional to the value of the parameter.

\end{document}
